% META: {"id": "TCOMPL-INEG-EX-006", "chapitre": "TCOMPL-INEGALITES", "type_objet": "exercice", "capacites": ["TCOMPL-INEG-C2"], "capacites_codes": ["C2"], "methodes": ["M2"], "parcours": 2, "competences": ["raisonner"], "duree_min": 15, "mode_creation": "ex_nihilo", "sources_inspiration": [], "fichier_tex": "chapitres/TCOMPL-INEGALITES/exercices/TCOMPL-INEG-EX-006.tex", "corrige_tex": "chapitres/TCOMPL-INEGALITES/corriges/TCOMPL-INEG-CO-006.tex", "status": "generated"}
% BEGIN-VERIFY
% from sympy import symbols, diff, factor, simplify, Rational
% x = symbols('x', nonnegative=True)
% L = x ** 3
% assert L.subs(x, 0) == 0 and L.subs(x, 1) == 1
% assert simplify(factor(x - L) - x * (1 - x) * (1 + x)) == 0
% for v in (Rational(1, 10), Rational(1, 2), Rational(9, 10), 1):
%     assert (x - L).subs(x, v) >= 0
% assert diff(L, x, 2) == 6 * x
% for v in (0, Rational(1, 2), 1):
%     assert diff(L, x, 2).subs(x, v) >= 0
% END-VERIFY
\begin{exercice}{TCOMPL-INEG-EX-006}{2}{15}
On modélise une courbe de Lorenz par $L(x)=x^3$ sur $[0\,;\,1]$.
\begin{enumerate}
  \item Vérifier que $L(0)=0$, $L(1)=1$ et $L(x)\leq x$ sur $[0\,;\,1]$.
  \item Montrer que $L$ est convexe sur $[0\,;\,1]$.
  \item Que signifierait $L(x)=x$ pour tout $x$ ?
\end{enumerate}
\end{exercice}
